\section{固定维度极小极大速率的证明}\label{app:minimax}

\begin{proof}
上下界采用相同的完整轨迹查询接口与期望轨迹预算。下界通过$\mathfrak C$中一个两点子模型构造：两模型的CPT梯度相距$\mathsf B^{-1/2}$阶，而停止后的查询记录在总变差距离下仍难以区分。上界则由定理~\ref{thm:estimator}中的无偏、有限方差且有限期望成本的输出给出。

\par\smallskip\noindent\textbf{第一步：可由CPT投资轨迹实现的单股票子模型。}\quad
考虑一只风险股票和现金，取$h=1$，交易成本为零，期初财富与参考点均为一。前期组合全部持有现金，可行性映射取恒等映射。风险股票的特征为$\R^d$中的第一标准基向量$e_1$，现金特征为零。在查询点$\theta$下，潜在风险资产权重满足
\[
 \mathsf U\sim\operatorname{Beta}(e^{\theta^{[1]}},1),
\]
其中$\theta^{[1]}$是参数的第一坐标。经式~\eqref{eq:action-map}的惯性映射后，执行的风险资产权重为$\chi\mathsf U$。

令风险收益$R=R_\star\xi$，其中$R_\star>0$固定且足够小，$\xi\sim\operatorname{Bernoulli}(p)$与$\mathsf U$独立。期末财富与参考点为
\[
 X_{\rm end}=1+\chi R_\star\mathsf U\xi,
 \qquad
 r_{\rm end}=1+\eta_+\chi R_\star\mathsf U\xi,
\]
故期末相对结果为
\begin{equation}\label{eq:minimax-relative-outcome}
 Y=\upsilon\mathsf U\xi,
 \qquad
 \upsilon=(1-\eta_+)\chi R_\star>0.
\end{equation}
这是定义~\ref{def:minimax-oracle}下可行的投资轨迹模型，其损失侧价值恒为零。

\par\smallskip\noindent\textbf{第二步：CPT梯度的局部分离。}\quad
令$v=v_+$，并定义
\[
 \mathfrak m(t)=\upsilon\,v'(\upsilon t),\qquad 0\le t\le1.
\]
平移幂函数满足$\mathfrak m(0)=0$，$\mathfrak m(t)>0$对$t\in(0,1]$成立，且$\mathfrak m$有界。由$\Pp(\mathsf U>t)=1-t^{e^{\theta^{[1]}}}$，在收益侧CPT积分中作变量替换$y=v(\upsilon t)$，得到
\begin{equation}\label{eq:minimax-cpt-submodel}
 J_p(\theta)
 =\int_0^1 \mathfrak m(t)
 w_+^\eps\!\left(p\{1-t^{e^{\theta^{[1]}}}\}\right)dt.
\end{equation}
在目标点$\theta_0=0$附近的紧邻域中，$t^{e^{\theta^{[1]}}}\abs{\log t}$在$(0,1)$上有可积包络，且$(w_+^\eps)'$有界，故可交换求导与积分。因此
\begin{equation}\label{eq:minimax-gradient-function}
 \nabla J_p(0)=\mathfrak g(p)e_1,
 \qquad
 \mathfrak g(p)
 =p\int_0^1\mathfrak m(t)(w_+^\eps)'(p(1-t))(-t\log t)\,dt.
\end{equation}
记
\[
 A_v=\int_0^1\mathfrak m(t)(-t\log t)\,dt>0,
 \qquad
 c_{w,0}=(w_+^\eps)'(0)>0.
\]
由于$\abs{(w_+^\eps)''}\le C_2$，对式~\eqref{eq:minimax-gradient-function}关于$p$求导，可得
\begin{align*}
 \mathfrak g'(p)
 &=\int_0^1\mathfrak m(t)(-t\log t)
 \left[(w_+^\eps)'(p(1-t))+p(1-t)(w_+^\eps)''(p(1-t))\right]dt\\
 &\ge (c_{w,0}-2C_2p)A_v.
\end{align*}
取充分小的$p_*>0$，使
\[
 \mathcal I=[p_*/2,3p_*/2]\subset(0,1/2],
 \qquad
 3C_2p_*\le c_{w,0}/2,
 \qquad
 c_*=c_{w,0}A_v/2.
\]
于是
\begin{equation}\label{eq:minimax-gradient-separation}
 \mathfrak g'(p)\ge c_*>0
 \qquad\text{对所有 }p\in\mathcal I.
\end{equation}

\par\smallskip\noindent\textbf{第三步：自适应停止查询记录中的信息量。}\quad
当$\mathsf B$足够大时，取
\[
 p_\pm=p_*\pm\delta,
 \qquad
 \delta=\delta_*/\sqrt{\mathsf B},
\]
并选择$\delta_*>0$使$p_\pm\in\mathcal I$。对固定内部区间$\mathcal I$上的伯努利分布，存在只依赖$\mathcal I$的有限常数$C_I$。由基本界$\operatorname{kl}(p\Vert q)\le (p-q)^2/\{q(1-q)\}$，
\begin{equation}\label{eq:bernoulli-kl-bound}
 \operatorname{kl}(p_+\Vert p_-)
 \le C_I\delta^2.
\end{equation}

考虑式~\eqref{eq:minimax-risk}中的任意可行估计器。分别以$\Pp_+^{\rm tr}$、$\Pp_-^{\rm tr}$表示两个模型下停止后的查询记录分布。每次查询位置可以依赖此前记录，但在给定任意查询点后，两模型的潜在Dirichlet动作分布完全相同；只有伯努利收益核不同。算法自身的随机化也可纳入查询记录，并在两模型之间保持相同。数据决定的查询次数$T_{\rm call}$是查询记录滤过下的停止时刻。停止似然比的链式法则给出
\begin{equation}\label{eq:stopped-transcript-kl}
 \operatorname{KL}(\Pp_+^{\rm tr}\Vert\Pp_-^{\rm tr})
 =\E_+T_{\rm call}\,\operatorname{kl}(p_+\Vert p_-)
 \le \mathsf B C_I\delta^2=C_I\delta_*^2.
\end{equation}
上式等号可通过在事件$\{T_{\rm call}\ge j\}$上逐次求和条件对数似然增量得到；该事件在第$j$次新收益观测前可测。由于$p_\pm$处于紧的内部区间且$\E_+T_{\rm call}<\infty$，这些增量一致可积。

选取足够小的$\delta_*$，使$C_I\delta_*^2\le1/8$。由Pinsker不等式，
\begin{equation}\label{eq:minimax-tv}
 \operatorname{TV}(\Pp_+^{\rm tr},\Pp_-^{\rm tr})\le\frac14.
\end{equation}

\par\smallskip\noindent\textbf{第四步：将估计问题化为检验问题。}\quad
由式~\eqref{eq:minimax-gradient-separation}及均值定理，两个梯度目标的距离满足
\begin{equation}\label{eq:minimax-target-distance}
 \operatorname{Sep}_{\mathsf B}
 :=\norm{\mathfrak g(p_+)e_1-\mathfrak g(p_-)e_1}_2
 \ge2c_*\delta.
\end{equation}
给定任意估计量$\widehat g$，以其更靠近的目标定义一次二元检验。若检验错误，则估计量距真实目标至少为$\operatorname{Sep}_{\mathsf B}/2$。两个检验错误概率之和至少为一减去总变差距离，故
\begin{align*}
 \max_{s\in\{+,-\}}
 \E_s\norm{\widehat g-\mathfrak g(p_s)e_1}_2^2
 &\ge \frac{(\operatorname{Sep}_{\mathsf B})^2}{8}
 \left\{1-\operatorname{TV}(\Pp_+^{\rm tr},\Pp_-^{\rm tr})\right\}\\
 &\ge \frac{3c_*^2\delta_*^2}{8\mathsf B}.
\end{align*}
两模型均属于$\mathfrak C$，对所有可行程序取下确界，即得到式~\eqref{eq:minimax-rate}的下界，其中$c_{\rm lb}=3c_*^2\delta_*^2/8$。

\par\smallskip\noindent\textbf{第五步：匹配预算的上界。}\quad
固定定理~\ref{thm:estimator}中任意可行的估计器设计$(n,p_{\rm act},b)$。定义~\ref{def:minimax-oracle}中的共同模型类常数使
\[
 \sup_{P\in\mathfrak C}\tr\Cov_P(Z_n^{\db})\le C_{\rm var}(n,p_{\rm act},b)<\infty,
 \qquad
 C_{\rm cost}=n(1+p_{\rm act} \Upsilon_b)<\infty
\]
对所有$P\in\mathfrak C$一致成立。取
\[
 M=\left\lfloor\frac{\mathsf B}{C_{\rm cost}}\right\rfloor,
\]
在$\theta_0$处独立运行$M$次未截断去偏估计，并求平均。其完整轨迹的总期望数为$MC_{\rm cost}\le\mathsf B$，故该程序可行。当$\mathsf B\ge2C_{\rm cost}$时，$M\ge \mathsf B/(2C_{\rm cost})$，因此
\[
 \sup_{P\in\mathfrak C}
 \E_P\norm{\widehat g_{M,n}^{\db}-g_P}_2^2
 \le\frac{C_{\rm var}(n,p_{\rm act},b)}{M}
 \le\frac{2C_{\rm var}(n,p_{\rm act},b)C_{\rm cost}}{\mathsf B}.
\]
于是式~\eqref{eq:minimax-rate}的上界成立，可取$C_{\rm ub}=2C_{\rm var}(n,p_{\rm act},b)C_{\rm cost}$及充分大的$\mathsf B_0$。上下界针对相同的目标、轨迹查询接口、损失函数和期望原始轨迹预算，定理得证。
\end{proof}
